\documentclass[11pt]{article}
\usepackage[a4paper, margin=1in]{geometry}
\usepackage{amsmath, amssymb, graphicx}
\usepackage{setspace}
\usepackage[numbers]{natbib}

\title{Research Statement}
\author{PremVijay Velmani}
\date{\today}

\begin{document}

\maketitle

\section*{Summary of My PhD Research}

My PhD research focuses on investigating the dynamical evolution of dark matter haloes and the impact of baryonic processes within them, using large-scale hydrodynamical simulations such as IllustrisTNG, EAGLE, and CAMELS. This work done under the supervision of Prof.Aseem Paranjape is relevant for the ongoing research in several aspects from cosmology and galactic astrophysics to particle physics of dark matter. More specifically, this helps to decouple the influence of baryonic astrophysical effects from the radial distribution of dark matter in these haloes. In addition to running and analyzing cosmological simulations, I have also worked with semi-analytic models of galaxy halo evolution.

\textbf{Key Contributions:}
\begin{itemize}
    \item \textbf{Characterization of Halo Relaxation:} Through statistical analysis of a large number of haloes in cosmological simulations IllustrisTNG and EAGLE, we have developed models of quasi-adiabatic relaxation that accurately predicts the change in the dark matter radial distribution in response to the galaxy formation and evolution. These results also revealed the significance of feedback-related effects on the relaxation response and explicit dependence on halo-centric distance. This model provides an accurate fit to the relaxation responses observed in simulations of dark matter haloes and assists in physical interpretation of the relaxation across a variety of haloes.
    
    \item \textbf{Feedback and Epoch Dependence:} Through extensive CAMELS simulations, We analyzed the effects of different feedback mechanisms on the relaxation at different cosmic times. My findings highlight that while the amount of feedback outflow strongly influences halo relaxation, its exact nature plays a lesser role.
    
    \item \textbf{Gas Equation of State and Feedback:} Using EAGLE simulations, We demonstrated that alongside feedback effects, the gas equation of state has a significant impact on the relaxation response of haloes. 

    \item \textbf{Dynamics of the relaxation response:}  We have studied the evolutionary history of the relaxation of haloes across cosmic time ($z=5$ to $z=0$). This revealed that in general the relaxation effects manifest immediately in the inner halo regions and with a time delay of around 2–3 Gyr in the outer halo regions.
    
    \item \textbf{Self-Similar Model for Halo Formation:} We developed a spherical self-similar model incorporating radiative cooling and artificial viscosity to simulate the co-evolution of gas and dark matter. This model offers a tractable method for understanding halo relaxation in simplified scenarios, consistent with results from full-scale simulations.
    
    \item \textbf{Implications for Baryonification Models:} The insights gained from these studies offer improvements to semi-analytical models and baryonification schemes by providing a better understanding of the coupling between baryonic and dark matter dynamics.
\end{itemize}

\section*{Postdoctoral Research Plans}

For my postdoctoral research, I plan to build on my PhD work by further probing the interaction between dark matter dynamics and baryonic processes, with the following goals:

\textbf{1. Extending the Quasi-Adiabatic Model:}  
I aim to extend the quasi-adiabatic model developed during my PhD to include more complex baryonic processes, such as star formation and black hole feedback, to better match observed galaxy properties.

\textbf{2. Investigating Halo Substructure Evolution:}  
I will study how baryonic processes affect subhalo dynamics and tidal stripping within massive dark matter haloes, particularly in the context of satellite galaxy evolution.

\textbf{3. Collaboration with Upcoming Surveys:}  
I will work on integrating simulation results with upcoming surveys such as LSST and Euclid, helping to bridge the gap between observational data and theoretical predictions of dark matter distribution and galaxy formation.

\textbf{4. Application to Baryonification Schemes:}  
Incorporating the knowledge gained from detailed simulations into efficient baryonification schemes for fast cosmological predictions.

\textbf{5. Exploring Alternative Dark Matter Models:}  
I also intend to explore how alternative dark matter models, such as self-interacting dark matter (SIDM), alter the relaxation properties of haloes and how these can be distinguished from standard cold dark matter (CDM) in both simulations and observations.

\textbf{Conclusion:}  
By combining my experience with large-scale cosmological simulations and theoretical modeling, I aim to advance our understanding of the role baryons play in shaping dark matter haloes and the observable Universe, while contributing to the ongoing development of semi-analytical models for next-generation galaxy surveys.

\end{document}
